\section{Назначение и область применения}

\subsection{Назначение программы}

Назначением программы является формирование синтетических данных аналитического
бенчмарка TPC-H \srcref{1} с заданным масштабом генерации и сохранением результатов в
форматах, пригодных для непосредственного использования в современных
аналитических системах и подсистемах хранения данных.

Программа предназначена для детерминированной генерации таблиц набора TPC-H,
их представления в колонoчной форме Apache Arrow \srcref{3} и записи результата в
поддерживаемые форматы хранения без промежуточного этапа ручной конвертации.
В отличие от традиционных генераторов, ограниченных выводом в
\texttt{tbl}/\texttt{csv} \srcref{4}, разрабатываемая программа обеспечивает
прямое формирование данных в форматах \texttt{Parquet} \srcref{5},
\texttt{ORC} \srcref{6}, \texttt{Lance} \srcref{8} и \texttt{Vortex}
\srcref{9} и допускает запись как на
локальную файловую систему, так и в \texttt{S3}-совместимые объектные
хранилища.

\subsection{Область применения программы}

Программа может применяться при разработке, испытаниях и исследовании
аналитических СУБД и OLAP-движков \srcref{2}, lakehouse-платформ,
подсистем хранения данных и
других программных средств, работающих с большими объемами колонoчных данных.

Основными областями применения программы являются:
\begin{enumerate}
    \item подготовка воспроизводимых наборов данных для функционального и
    нагрузочного тестирования аналитических систем;
    \item исследование характеристик различных форматов хранения и стратегий
    записи данных;
    \item проверка корректности работы систем чтения, сериализации и обработки
    колонoчных данных;
    \item проведение сравнительных экспериментов, связанных с локальным и
    удаленным размещением аналитических наборов данных.
\end{enumerate}

Программа ориентирована на использование разработчиками, исследователями,
инженерами по тестированию и иными специалистами, которым требуется быстрое,
детерминированное и масштабируемое получение данных набора TPC-H в форме,
непосредственно пригодной для дальнейшего анализа и эксплуатации.
